% Authoritative LaTeX revision; do not regenerate from Markdown.
\chapter{Combinatorial search and complexity}
\Callout{Revision boundary}{This chapter revises the standalone review candidate. The previously accepted Chapters 1--2 and their source hashes remain unchanged. This combined volume is a new review candidate, not a new acceptance record.}
\section{Before scaling the experiment, name its resource}
The preceding chapter followed molecular candidates through selective
operations. Before making that experiment larger, pause at a simpler question:
what exactly became faster when candidates were processed together?

Consider four candidates and three dependent predicate tests. In an abstract
machine with enough parallel test slots, all four candidates can undergo the
first test together, then the second, then the third. The elapsed number of
dependent stages is three, but twelve candidate-tests have still occurred.
The illustration counts every test, including tests on candidates that might
already have failed; early stopping would be a different work policy.

\NativeFigure{work-depth}{Rows are candidates; columns are dependent test
stages. Horizontal arrows mean precedence for the same candidate. A column
can execute in parallel only under the declared resource model. Twelve tests
and depth three are different counts, neither a measured reaction time.}{fig:textbook-work-depth}

Let $C$ be candidate count and $S$ the number of sequential stages. Under this
idealized all-candidates/all-stages policy,
\begin{equation}
W=CS\quad[\text{candidate-tests}],\qquad
D=S\quad[\text{dependent stages}].
\end{equation}
The dimensions expose the mistake in equating work with time. Converting
$D$ into seconds requires a duration model for a stage. Implementing the
parallel columns requires capacity for the candidates and the operations.
Neither requirement disappears when a diagram draws the column vertically.

\textbf{Prediction problem.} Suppose candidate count doubles while stage count
stays fixed. What happens to these two counts? Does that prove the physical
experiment takes the same time?

\textbf{Worked answer.} Work doubles from 12 to 24 candidate-tests and abstract
depth remains three. Equal physical time does not follow: concentration,
mixing, reaction completion, separation, transfers and readout may change.
The formula is a declared abstract model. The existing companion function
\texttt{work\_depth} computes its counts; it does not simulate chemistry.

There is a second, independent scaling question. A population can have many
copies without representing every relevant candidate. If one independently
drawn copy produces a detectable witness with probability $q$, then $m$ copies
miss the witness with probability $(1-q)^m$. Thus the probability of observing
at least one witness is
\begin{equation}
P_{\mathrm{signal}}=1-(1-q)^m.
\end{equation}
This calculation assumes independent copies and a fixed per-copy probability.
It is not a calibrated molecular yield model. Correlated synthesis failures
or concentration-dependent reactions can invalidate it.

The rest of the chapter keeps these distinctions separate: the logical
decision problem, the algorithmic work/depth model and the physical evidence
needed to justify a conclusion. A short sequence of stages alone cannot prove
that an exponentially large candidate population is cheap to construct.



\section{A successful experiment leaves an algorithmic
question}\label{a-successful-experiment-leaves-an-algorithmic-question}

Chapter 2 showed how a molecular population can encode candidates and
how selective operations can enrich a witness. But a seven-vertex
success cannot, by itself, answer a question about a growing family of
inputs. What happens as the graph grows? Which resource grows? Does an
observed failure mean the graph has no solution, or only that the
experiment failed to expose one?

To ask those questions precisely, separate an \textbf{instance}, an
\textbf{algorithm}, and an \textbf{implementation}. The instance is a
particular directed graph with specified endpoints. The algorithm
determines which transformations and tests are performed. The
implementation supplies storage, arithmetic, reactions, transfers and
measurements. Complexity theory begins with an explicit abstract cost
model. Molecular engineering must additionally explain how that model is
realized.

We will use a new four-vertex graph, not silently replace the historical
graph:

\begin{equation}
V=\{0,1,2,3\},\quad s=0,\quad t=3,
\end{equation} \begin{equation}
E=\{(0,1),(0,2),(1,2),(2,1),(1,3),(2,3)\}.
\end{equation}

It has two Hamiltonian witnesses: \((0,1,2,3)\) and \((0,2,1,3)\). The
existence of two paths will matter: an algorithm that remembers only one
representative can still answer an existence question correctly, while
losing counting information.

\section{The quantifier makes the
problem}\label{the-quantifier-makes-the-problem}

Define a relation \(R(G,s,t,P)\) that holds exactly when \(P\) lists
every vertex once, starts at \(s\), ends at \(t\), and follows only
edges of \(G\). The decision question is

\begin{equation}
\exists P\;R(G,s,t,P)?
\end{equation}

The search question asks for such a \(P\), or a justified report that
none exists. The verification question takes a proposed \(P\) as an
additional input and evaluates \(R\). These questions share a relation
but require different outputs. Rejecting \((0,1,3,2)\) does not answer
the existence question for our graph; both valid paths above still
exist.

\DeepFigure{DNAD-03-F1}{\textbf{DNAD-03-F1.} A verified certificate establishes a yes-instance.
A rejected certificate only eliminates the submitted order. The
distinction is the existential quantifier, not a difference in
terminology.}{fig:DNAD-03-F1}

Under an adjacency-matrix representation already in memory, a verifier
can mark each visited vertex and check each successive edge. This takes
\(O(n)\) word-level operations for an \(n\)-vertex candidate when
indexing and labels fit the chosen word model. Reading or constructing
the full matrix still costs \(\Theta(n^2)\) bits/entries at the
appropriate representation level. Do not quote the verifier's inner loop
as the end-to-end cost of acquiring the graph.

The reference \texttt{verify} validates its graph argument before
checking the certificate, so its total Python work also includes graph
normalization. Set lookup, integer representation and interpreter
overhead are implementation details; our operation counts below are not
measured processor timings.

\section{Polynomial in which size?}\label{polynomial-in-which-size}

An algorithm's running time is a function of the encoded input length,
not merely the largest number printed in its input. A binary integer
\(N\) takes approximately \(\log_2 N\) bits; performing \(N\) iterations
can therefore take exponentially many iterations in that bit length. For
explicitly represented simple graphs, matrix and adjacency-list
encodings differ in size but can be translated with polynomial overhead.
A succinct circuit describing an exponentially larger graph would define
a different representation problem.

The usual class \textbf{P} contains decision problems decidable in
polynomial time. \textbf{NP} can be characterized by polynomial-length
certificates whose validity is checked in polynomial time. The
certificate definition does not prescribe enumerating all certificates,
and it does not describe a physical machine creating every possibility
for free. These are the conventions used in
\href{https://ocw.mit.edu/courses/18-404j-theory-of-computation-fall-2020/45e2fd621349cfd7c9faf93a6ba134a3_MIT18_404f20_lec14.pdf}{Sipser's
lecture on P, NP and reducibility}.

An important negative statement follows from logic alone: the definition
supplies short certificates for yes-instances. It does not automatically
supply short certificates for no-instances. Negating the entire
existence statement gives \(\forall P\,\neg R(G,s,t,P)\), not the
rejection of one chosen \(P\). We should not confuse that distinction
with a proof that no concise negative certificate could ever exist.

\DeepFigure{DNAD-03-F5}{\textbf{DNAD-03-F5.} Our four-vertex matrix contains 16 adjacency bits,
not 16 bits for the entire instance: the representation must also supply
its size and endpoints. The integer 1024 needs 11 binary digits. Six
physical copies of one sequence are six molecules but only one distinct
sequence design. None of these counts can substitute for another.}{fig:DNAD-03-F5}

Here is a precise way to read the integer example. For positive \(N\),
its ordinary binary length is \(\lfloor\log_2 N\rfloor+1\). Setting
\(N=2^k\) therefore makes a loop of \(N\) iterations exponential in the
\(k+1\) input bits. By contrast, explicitly listing \(N\) objects
already consumes at least \(N\) representation units. The same loop
bound can have a different complexity interpretation because the input
representation changed.

The class \textbf{coNP} is defined by complementing languages:
\(L\in\mathrm{coNP}\) exactly when \(\overline L\in\mathrm{NP}\). It is
not the collection of all problems outside NP. These conventions are
stated in
\href{https://www.cs.toronto.edu/~toni/Courses/Complexity2015/lectures/lecture3.pdf}{Toronto's
complexity lecture, section 4}.

We can derive the inclusion \(P\subseteq NP\cap coNP\) without guessing
which classes are equal. A deterministic polynomial-time decider is also
a verifier that ignores its certificate, so \(P\subseteq NP\). Reversing
the output of a deterministic decider still decides its complement in
polynomial time. Thus that complement belongs to NP too, placing the
original language in coNP. Reversing the answer on just one
nondeterministic branch does not perform this operation on the whole
existence statement.

\DeepFigure{DNAD-03-F7}{\textbf{DNAD-03-F7.} The four candidate slots are a schematic finite
universe, not the exhaustive orders of our running graph. The left panel
needs one accepted candidate; the right panel needs all candidates
rejected. Checking only the four drawn slots would not establish a
negative answer for a problem with further candidates. The class
definitions below the panels do not assert strict inclusions.}{fig:DNAD-03-F7}

\section{Reuse the future question, not the entire
history}\label{reuse-the-future-question-not-the-entire-history}

Naive fixed-endpoint enumeration examines up to \((n-2)!\) interior
orders. Many prefixes, however, lead to the same remaining problem.
Suppose two valid prefixes have visited the same set \(S\) and both end
at vertex \(v\). The remaining vertices are \(V\setminus S\), and the
next step must use an outgoing edge from \(v\). Their internal orders no
longer affect whether an unused suffix can complete the path.

This motivates a Boolean state:

\begin{equation}
D[S,v]=\text{a simple path from }s\text{ visits exactly }S\text{ and ends at }v.
\end{equation}

\DeepFigure{DNAD-03-F2}{\textbf{DNAD-03-F2.} This separate five-vertex illustration exposes a
state equivalence before the final step. Equal visited sets alone are
insufficient: different last vertices can have different outgoing edges.
Equal last vertices alone are also insufficient: different used vertices
permit different future moves.}{fig:DNAD-03-F2}

Start with \(D[\{s\},s]=\mathrm{true}\). For a reachable state and an
unused neighbor \(w\), set

\begin{equation}
D[S\cup\{w\},w]\leftarrow\mathrm{true}
\quad\text{if }(v,w)\in E\text{ and }w\notin S.
\end{equation}

Process sets in increasing cardinality. The invariant is proved by
induction. The base state is the one-vertex path. Every transition
appends one unused vertex through a permitted edge, preserving validity.
Conversely, remove the last vertex from any valid longer path. Its
remaining prefix is a valid smaller state, and the recurrence restores
exactly that final step. Therefore \(D[V,t]\) is true if and only if a
witness exists.

Our implementation additionally refuses to enter \(t\) before the set is
full. This is safe because a simple path ending at \(t\) cannot leave
\(t\) and later revisit it. One predecessor per reachable state is
enough to reconstruct one witness. It is not enough to count all
witnesses or represent their molecular multiplicities. Different
questions need different state values and sometimes different state
definitions.

There are at most \(n2^n\) mask/endpoint pairs. Scanning up to \(n\)
outgoing neighbors for each gives an \(O(n^2 2^n)\) upper bound in the
usual unit-cost subset-DP account, with \(O(n2^n)\) state storage. This
is not a best-known-algorithm claim. The executable implementation uses
arbitrary-precision Python bitmasks, whose operations acquire a
size-dependent cost as masks exceed machine words. Both accounts remain
exponential, not polynomial.

\section{What the counters actually
show}\label{what-the-counters-actually-show}

On the four-vertex example, the implementation creates six reachable
states, arranged in cardinality layers \(1,2,2,1\), and performs ten
outgoing-neighbor scans. The two full paths merge into one terminal
reachability state. These counts come from the executable result
artifact (companion file: results.json), not a decorated complexity
diagram.

For the complete directed graph with fixed endpoints and \(n\ge3\), our
early-terminal rule permits exactly

\begin{equation}
2+(n-2)2^{n-3}
\end{equation}

reachable states. To derive this, count the initial state and final
state separately. Every other state chooses one of the \(n-2\) interior
vertices as its last vertex, then any subset of the other \(n-3\)
interior vertices as already visited. Completeness of the graph makes
every such choice reachable. At \(n=8\), the result is 194 states. This
is a state-count formula for this graph family and this implementation,
not a general physical-resource law.

The tests compare subset search against independent permutation
enumeration on \textbf{every one of the 4,096 simple directed
four-vertex graphs}. This is useful fault detection. The inductive
invariant is what supports correctness for arbitrary finite graph size;
exhaustive testing at size four cannot replace it.

\DeepFigure{DNAD-03-F8}{\textbf{DNAD-03-F8.} The curve uses complete directed graphs, unlike our
six-edge running graph. At \(n=4\), the complete graph has 15 neighbor
scans while the running graph has ten. At \(n=8\), there are 720
possible interior orders, 194 reachable states and 1,351 scans. These
are generated counts, not measured execution times.}{fig:DNAD-03-F8}

Why can scans outnumber candidate orders at this small size? A neighbor
scan is an attempted local extension, including extensions rejected
because the neighbor was already visited or is the terminal vertex too
early. An enumerated order is a whole proposed route, which itself needs
edge checks. Equating one scan with one verified order compares
different units. Also, a decision-only enumerator can stop at its first
witness; our permutation oracle deliberately lists every witness. The
comparison describes full enumeration and state exploration, not a fair
early-stopping runtime contest.

The asymptotic state reduction is nevertheless meaningful. With
\(k=n-2\) interior vertices, full enumeration has \(k!\) orders while
the dense reachable-state count is \(2+k2^{k-1}\). The ratio of
successive factorial counts is \(k+1\), while the state count eventually
grows by a factor close to two. This explains why merging histories
matters at increasing sizes, without turning an exponential method into
a polynomial one. Memory, transition work and reconstruction costs still
require their own accounting.

\section{A reduction has a direction and two
implications}\label{a-reduction-has-a-direction-and-two-implications}

A polynomial many-one reduction from problem \(A\) to problem \(B\)
computes a mapping \(f\) such that

\begin{equation}
x\in A\iff f(x)\in B.
\end{equation}

Thus a fast solver for \(B\), composed with \(f\), would solve \(A\).
Hardness travels from the source problem toward the target problem. A
mapping in the reverse direction establishes a different claim.

Consider a directed Hamiltonian-cycle instance. Choose a pivot vertex,
replace it by a source \(s\) and sink \(t\), redirect outgoing pivot
edges from \(s\), and redirect incoming pivot edges into \(t\). Keep the
other edges. A cycle through the pivot opens into a spanning
\(s\)-to-\(t\) path. Conversely, identifying the endpoints of such a
path closes a cycle through the original vertices. The construction adds
one vertex and preserves the number of edges for our loop-free graph
convention. This is the vertex-splitting reduction presented in
\href{https://ocw.mit.edu/courses/6-046j-design-and-analysis-of-algorithms-spring-2015/1dee182d301235b901119a0ff57f05e2_MIT6_046JS15_Recitation8.pdf}{MIT's
algorithms recitation}.

\DeepFigure{DNAD-03-F3}{\textbf{DNAD-03-F3.} The three-cycle is an illustration, not the proof
for every graph. The code checks all simple directed three-vertex graphs
and every pivot, including graphs without cycles.}{fig:DNAD-03-F3}

This argument transfers an established hardness premise for directed
Hamiltonian cycle to the fixed-endpoint path problem. It does not
establish that premise from scratch. A full path from satisfiability to
the Hamiltonian family needs a separately reviewed construction; do not
hide that missing step inside a decorative arrow.

\section{Encode the whole predicate in Boolean
clauses}\label{encode-the-whole-predicate-in-boolean-clauses}

We can also map a fixed-endpoint path instance to a satisfiability
instance. This direction is useful for expressing constraints but, by
itself, is not a proof that the path problem is NP-hard.

Let \(x_{v,p}\) mean that vertex \(v\) occupies position \(p\). There
are \(n^2\) variables. Require exactly one vertex at each position and
exactly one position for each vertex. An elementary exactly-one encoding
uses one positive clause for ``at least one'' and a negative two-literal
clause for every pair that must not both be true. Add the unit clauses
\(x_{s,0}\) and \(x_{t,n-1}\). Finally, for each forbidden transition
\((u,v)\notin E\) and each \(p<n-1\), add

\begin{equation}
\neg x_{u,p}\lor\neg x_{v,p+1}.
\end{equation}

\DeepFigure{DNAD-03-F4}{\textbf{DNAD-03-F4.} A permutation matrix can still describe a
forbidden-edge route. The transition clauses implement the extra
adjacency condition already exposed by Chapter 2's pseudo-path.}{fig:DNAD-03-F4}

Soundness is direct: every satisfying assignment is a permutation
matrix, and the endpoint and transition clauses turn its decoded order
into a witness. Completeness is also direct: encode any valid path by
putting true in its occupied vertex/position cells; every clause family
is then satisfied.

Our encoding emits

\begin{equation}
2n+2n\binom n2+2+(n-1)(n^2-m)
\end{equation}

clauses for a loop-free directed graph with \(m\) edges. The four terms
count positive exactly-one clauses, pairwise exclusions, endpoints and
forbidden transitions. Self-transitions are included among forbidden
edges even though the permutation constraints already rule them out;
redundancy does not invalidate the encoding. Our toy graph produces 16
variables, 88 clauses and 190 literal occurrences. Tests inspect all
\(2^9\) Boolean assignments for each three-vertex graph, not just
hand-picked permutation assignments. That catches encodings that
accidentally admit malformed matrices.

\section{Turn a decision oracle into a
witness}\label{turn-a-decision-oracle-into-a-witness}

Suppose an oracle answers the fixed-endpoint existence question. Ask it
once about the original graph. If the answer is yes, consider edges one
by one. Delete an edge whenever the oracle says a witness still exists
without it. This maintains a yes-instance and terminates after at most
\(m+1\) oracle calls.

Why does the remaining graph reveal a path? Select any spanning witness
in the final graph. If an edge lay outside that witness, removing it
would preserve that witness, contradicting edge minimality. Hence every
remaining edge belongs to the chosen path; the final graph contains
exactly its \(n-1\) edges. An edge that was indispensable when
considered cannot become dispensable after further deletions: deleting
more edges cannot create a missing path. This justifies a single pass
over the edges.

The companion implements this reduction with \texttt{subset\_search} as
its oracle. That oracle is exponential; calling it ``an oracle'' does
not remove its cost. On our six-edge toy graph, seven calls leave the
witness \((0,2,1,3)\). This differs from the first witness returned by
subset search, without any contradiction: the task requests a witness,
not a unique canonical order.

\DeepFigure{DNAD-03-F6}{\textbf{DNAD-03-F6.} A NO answer describes the trial graph without the
tested edge. It does not describe the retained graph: on NO, the
algorithm keeps that edge and preserves a yes-instance. This distinction
is visible at query 3, where deleting \(0\to2\) would remove the only
surviving way to leave the source after \(0\to1\) has already been
deleted.}{fig:DNAD-03-F6}

We can sharpen the edge-minimality argument by naming the time
dependence. Let \(G_i\) be the retained graph before testing edge
\(e_i\), and let \(G_f\subseteq G_i\) be the final retained graph. If
\(G_i-e_i\) has no witness, neither can its subgraph \(G_f-e_i\).
Therefore a retained edge never needs to be reconsidered. If the final
graph had an extra edge outside a spanning witness \(P\), that edge
would contradict this fact because \(P\) would survive its deletion.
This proves both the one-pass rule and the final \(n-1\)-edge structure.

The trace records the graph \emph{after} each decision, rather than just
a list of answers. Tests independently enumerate witnesses in both the
trial and retained graph at each step. This catches a subtle
implementation failure: printing correct oracle answers while
accidentally retaining the wrong edge set.

\section{Bring the accounting back to
molecules}\label{bring-the-accounting-back-to-molecules}

The Boolean DP can forget alternate histories because its question is
existential. A molecular population cannot always be treated that way.
Copy numbers affect survival and detection, and different sequences with
the same abstract future role may have different physical interactions.
A proposed molecular implementation of the recurrence must explain how
sets and endpoints are represented, how equivalent states are
recognized, and what each transition costs.

Similarly, reducing a graph to a polynomial-size formula says that the
formula is not exponentially larger than the encoded graph. It does not
say that solving the formula is polynomial. Putting candidate
assignments into many simultaneous physical objects changes allocation
and depth, not the need to account for those objects.

The next resource chapter should therefore keep total work,
critical-path depth, material, memory and readout separate. No theorem
in this draft turns all DNA Hamiltonian solvers into factorial
algorithms, and no small-instance simulation proves asymptotic
superiority over electronic computation.

\section{Parallel depth is not total
work}\label{parallel-depth-is-not-total-work}

A parallel algorithm has a dependency structure. If four independent
candidates each pass through three successive abstract tests, there are
twelve test applications but only three tests on any one candidate's
critical path. With enough workers, the tests at each level can overlap.
That does not remove the other nine applications; it changes when and
where they are carried out.

\DeepFigure{DNAD-03-F9}{\textbf{DNAD-03-F9.} Every circle is one abstract test. Columns are
dependency levels, not laboratory time points. Work is the total number
of circles; depth is the length of the longest dependent chain. No
reaction-rate calibration is implied.}{fig:DNAD-03-F9}

Write total work as W, dependency depth as D, and available processors
as P. Under an ideal discrete model in which each operation occupies one
processor for one time unit, execution time satisfies

\begin{equation}
T_P\ge\max\{W/P,D\}.
\end{equation}

The work bound follows because P processors complete at most P
operations per unit; the depth bound follows because a dependent
operation cannot precede its input. These are lower bounds, not a
scheduling guarantee. Communication, allocation, load imbalance and
synchronization can increase time further. The fixture in
\texttt{completion\_diagnostics.py} gives W=12 and D=3. Four ideal
workers can realize the three levels, while one worker needs at least
twelve units.

A test tube is not a collection of ideal processors. Molecules interact
at concentration-dependent rates, may compete for substrates and may be
lost during handling. Thus P cannot simply be replaced by molecule count
to predict a reaction duration. The transferable lesson is the
accounting separation: a short dependency chain does not imply a small
allocation, small total work or cheap readout. Chapter 4 will add those
physical coordinates rather than rename depth as molecular speed.

\section{Missing signal has several
explanations}\label{missing-signal-has-several-explanations}

There are at least three distinct ways a positive instance can produce
no recorded witness: no witness copy enters the sample; a copy enters
but does not survive the operations; a surviving copy does not produce a
recorded signal. These alternatives concern physical completeness, not
logical soundness.

\DeepFigure{DNAD-03-F10}{\textbf{DNAD-03-F10.} The panels are alternative explanations, not a
serial protocol. Curved marks denote schematic molecules, not counted
microscopic observations. The probability formula below them is an
explicit toy model, not a fitted assay.}{fig:DNAD-03-F10}

Suppose M independent candidate copies are drawn. Let p be the
probability that one is a witness, s the conditional probability that a
witness survives, and d the conditional probability that a surviving
witness is detected. The chain rule gives a per-copy recorded-witness
probability q=psd; this factorization uses conditional probabilities and
does not require the three events within a copy to be independent. The
additional assumption is independence \textbf{across} copies. Only then
is

\begin{equation}
\Pr(\text{at least one recorded witness})=1-(1-psd)^M.
\end{equation}

For p=.1, s=.5 and d=.8, q=.04. Ten copies give about .335167
probability of a recorded witness, so no signal remains quite plausible
even when a witness is possible. The implementation uses \texttt{log1p}
and \texttt{expm1} to avoid avoidable cancellation for small q, and
handles zero-copy and certain-signal boundaries explicitly. Tests
compare it with direct binomial reasoning on small cases.

Shared contamination, common reagent failure and competition can
correlate copies. In that case the exponent formula need not hold. A
positive control can show that some detection path functioned, but
cannot prove all witness species were represented or retained.
Conversely, this toy model cannot turn a signal into a valid path: that
still requires decoding and the independent verifier.

\section{What is proved, what is imported, what is
measured}\label{what-is-proved-what-is-imported-what-is-measured}

\DeepFigure{DNAD-03-F11}{\textbf{DNAD-03-F11.} Opening and closing a cycle are inverse witness
constructions under the stated loop-free, distinct-endpoint conventions.
The drawing is one example; the argument and encoding-size account
justify the general reduction.}{fig:DNAD-03-F11}

This chapter proves its subset recurrence, path-to-CNF equivalence and
decision-to-search construction. It proves how directed cycle hardness
would transfer to the fixed-endpoint path problem. The upstream
NP-hardness result is an attributed theoretical premise, not a proof
reconstructed from the original Cook and Karp papers here. Stating that
boundary is more useful than presenting an incomplete satisfiability
gadget as if its local arrows established a theorem.

Keep three evidence columns separate. A proof establishes a quantified
statement under assumptions. Finite tests challenge an implementation of
that statement. A laboratory measurement constrains a physical system
and its measurement model. None can silently replace either of the
others. A theorem does not certify a pipetting operation, and a
successful reaction on one graph does not prove a complexity-class
separation.

\section{The physical bridge the abstraction must
preserve}\label{the-physical-bridge-the-abstraction-must-preserve}

\DeepFigure{DNAD-03-F12}{\textbf{DNAD-03-F12.} A duplex illustrates that directed sequences have
a physical realization, but the DP table is not an assay. One true
Boolean entry does not distinguish one remaining molecule from many. The
cartoon's four bases illustrate orientation; they are not a proposed
stable reagent or a molecular DP design.}{fig:DNAD-03-F12}

When two histories merge into D{[}S,v{]}, the algorithm preserves the
future existence question. A proposed molecular merger must additionally
preserve whatever physical information later operations require.
Sequence identity, copy number and accessibility may affect those
operations even if the abstract visited set and endpoint match.
Therefore the implementer must provide a representation map and an
error/resource contract, not just rename a molecule as a state. This is
the precise handoff to molecular parallelism and resource accounting in
Chapter 4.

\section{Exercises with worked
reasoning}\label{exercises-with-worked-reasoning}

\begin{enumerate}
\def\labelenumi{\arabic{enumi}.}
\item
  \textbf{A failed certificate.} In the four-vertex example, submit
  \((0,1,3,2)\). Does rejection prove the instance is negative?
  \textbf{Solution:} It ends at 2 rather than 3 and requires absent edge
  \(3\to2\). This rejects that order only; either listed witness proves
  the instance positive.
\item
  \textbf{Insufficient state.} Why not store only the visited set?
  \textbf{Solution:} Two prefixes using the same set may end at
  different vertices. Their available outgoing edges differ, so they
  need not admit the same completion. The last vertex must be retained
  unless another representation supplies equivalent information.
\item
  \textbf{Counting witnesses.} Can the one-predecessor table count both
  toy witnesses? \textbf{Solution:} No.~Boolean reachability merges
  successful histories. Counting requires a sum over predecessor counts
  with a base count of one. Reconstructing all paths requires still more
  information or traversal; its output can itself be large.
\item
  \textbf{Reduction direction.} Does a polynomial path-to-SAT encoding
  prove path hardness? \textbf{Solution:} Not by itself. It says a SAT
  solver can solve the encoded path problem. For a hardness transfer
  from SAT, a reduction must go from SAT to the claimed hard target,
  with both existence implications proved.
\item
  \textbf{Clause audit.} Explain the 88 toy clauses. \textbf{Solution:}
  Eight positive row/column clauses plus 48 pairwise exclusions plus two
  endpoint units plus \(3(16-6)=30\) transition exclusions total 88. The
  redundant self-transition exclusions are included in that total.
\item
  \textbf{Oracle budget.} Why do seven oracle calls not constitute a
  seven-step polynomial-time path solver? \textbf{Solution:} The calls
  contain substantial computations. Here each call uses an exponential
  subset algorithm. A polynomial number of calls to a hypothetical
  polynomial-time oracle would imply a polynomial-time composition, but
  that hypothesis has not been established.
\item
  \textbf{Input-size trap.} A routine takes \(N^2\) iterations on a
  binary integer \(N=2^k\). Is it polynomial-time in the input length?
  \textbf{Solution:} Its input has \(k+1\) bits, while the iteration
  count is \(2^{2k}\). A polynomial expression in the numeric value is
  not a polynomial in the encoded length.
\item
  \textbf{Complement closure.} Why is a problem in P also in coNP?
  \textbf{Solution:} Flip the output of its deterministic
  polynomial-time decider to decide the complement. That complement is
  in P and therefore NP, which is exactly the required coNP condition.
  No claim that NP is closed under complementation is needed.
\item
  \textbf{Audit the plotted graph.} Why do the two four-vertex examples
  report ten and fifteen scans? \textbf{Solution:} The running graph has
  six edges, whereas the plotted complete graph has twelve. The
  reachable-state count is six in both, but scanning outgoing adjacency
  at those states incurs different work. State count alone does not fix
  transition work.
\item
  \textbf{Trace an invariant.} After deleting \(0\to1\), why retain
  \(0\to2\)? \textbf{Solution:} Removing the latter leaves no outgoing
  edge from source 0 and hence no spanning source-to-target path.
  Retaining it preserves the path \(0,2,1,3\). The oracle's NO refers to
  the trial deletion, not to the retained graph.
\item
  \textbf{Design a differential test.} How would you test a trace
  without trusting the DP oracle? \textbf{Solution:} For every small
  graph, enumerate all interior permutations independently. Replay each
  proposed deletion, compare its existence answer with enumeration, and
  check that the retained graph still has a witness. Verify that the
  final edge set is exactly the returned path's consecutive pairs.
  Passing finite tests supplements, rather than replaces, the
  monotonicity proof.
\item
  \textbf{A representation audit.} An experiment advertises ``only \(n\)
  DNA designs,'' with \(2^n\) copies of each. What has the headline
  omitted? \textbf{Solution:} It has counted distinct designs but
  omitted the \(n2^n\) physical copies stipulated by its own inventory.
  Also request strand lengths, allocation across designs, transfer
  losses and detection assumptions. This is an audit of that proposal,
  not a lower bound on every molecular algorithm.
\end{enumerate}

\section{Publication work still
required}\label{publication-work-still-required}

All twelve storyboard figures are now produced with semantic TXT
companions. The standalone Chapter 3 review edition includes a source
list, glossary and tested code appendix. It does not modify the accepted
cumulative edition. Independent subject review and cumulative LaTeX
index/bibliography integration remain release gates; the imported
upstream hardness theorem is not claimed as a newly reconstructed full
proof.

\section{Selected glossary}\label{selected-glossary}

\textbf{Certificate:} a proposed finite witness checked by a verifier.

\textbf{Reduction:} an efficiently computable map preserving a decision
predicate.

\textbf{Work:} total elementary operations in the declared model.

\textbf{Depth:} longest chain of dependent operations in that model.

\textbf{Physical completeness:} conditions under which a logically
existing witness is represented, retained and observed; not implied by
predicate soundness.

\section{Appendix A. Executable boundary
diagnostics}\label{appendix-a.-executable-boundary-diagnostics}

These functions are reproduced from completion\_diagnostics.py and
exercised by the repository tests. The full reference program,
exhaustive checks and machine-readable results remain alongside the
manuscript in drafts/ch03/. These are teaching calculations, not
calibrated biological measurements or a trained model.

\begin{verbatim}
import math

def coverage(copies, probability, survival=1., detection=1.):
    if type(copies) is not int or copies < 0:
        raise ValueError("copies must be a nonnegative integer")
    for p in (probability, survival, detection):
        if not math.isfinite(p) or not 0 <= p <= 1:
            raise ValueError("probabilities must lie in [0,1]")
    q=probability*survival*detection
    if copies == 0 or q == 0: return 0.
    if q == 1: return 1.
    return -math.expm1(copies*math.log1p(-q))

def work_depth(candidates, stages):
    if any(type(v) is not int or v < 1 for v in (candidates,stages)):
        raise ValueError("positive integer dimensions required")
    return {"candidates":candidates,"stages":stages,
            "work":candidates*stages,"depth":stages}
\end{verbatim}

\section{Appendix B. Sources and evidence
boundaries}\label{appendix-b.-sources-and-evidence-boundaries}

\subsection{Chapter 3 source ledger}\label{chapter-3-source-ledger}

\subsubsection{Review-candidate extension, 12 September
2026}\label{review-candidate-extension-12-september-2026}

The earlier entries below retain the access record of the first two
passes. The new work/depth and missing-signal examples are original
derivations with explicit assumptions, not laboratory measurements. For
the signal calculation, survival and detection are conditional
probabilities in a per-copy chain; independence is assumed across
candidate copies only when exponentiating the no-signal probability. The
implementation is checked against explicit two-copy enumeration,
boundary cases and a small-probability stability case.

MIT 6.046J Recitation 8 was rechecked for the split-vertex reduction.
This chapter proves that map and explicitly imports the hardness of
directed Hamiltonian cycle as a premise. It does not claim to supply the
entire Cook-Karp hardness chain. Figures F9-F12 separate abstract work,
molecular population, observation and proof obligations. All twelve
figures are produced; calibrated physical accounting and independent
scientific review remain open.

Consulted 9 September 2026. This ledger distinguishes directly read
support from located-but-unread material. It is not a systematic review
or independent approval.

\subsubsection{Directly read support}\label{directly-read-support}

\begin{itemize}
\tightlist
\item
  Michael Sipser, MIT 18.404J (2020),
  \href{https://ocw.mit.edu/courses/18-404j-theory-of-computation-fall-2020/45e2fd621349cfd7c9faf93a6ba134a3_MIT18_404f20_lec14.pdf}{Lecture
  14: P and NP, SAT, polynomial-time reducibility}. The certificate,
  input-encoding and reduction discussion was read from the instructor's
  notes. Used for class conventions, not a proof of a separation between
  P and NP.
\item
  MIT 6.046J (2015),
  \href{https://ocw.mit.edu/courses/6-046j-design-and-analysis-of-algorithms-spring-2015/1dee182d301235b901119a0ff57f05e2_MIT6_046JS15_Recitation8.pdf}{Recitation
  8}, section 2, physical page 2. Read the directed-cycle-to-path
  vertex-splitting construction and both existence implications. Our
  original figure uses a three-cycle; the executable extension exposes
  fixed endpoints and arbitrary pivot relabeling. The source's naive
  verifier bounds are not substituted for our separately stated
  representation-aware account.
\item
  Sipser,
  \href{https://ocw.mit.edu/courses/18-404j-theory-of-computation-fall-2020/resources/lecture-15-np-completeness/}{Lecture
  15 resource and transcript}. Read the caution that directed
  3SAT-to-HAMPATH gadgets require internal consistency details. No
  incomplete gadget diagram was promoted to a complete proof in this
  draft.
\end{itemize}

\subsubsection{Located but not used as a fully read
proof}\label{located-but-not-used-as-a-fully-read-proof}

\begin{itemize}
\tightlist
\item
  Stephen Cook, \emph{The Complexity of Theorem-Proving Procedures}
  (1971):
  \href{https://www.cs.toronto.edu/~sacook/homepage/1971.pdf}{author-hosted
  scan}, linked from the author's publication list. The bibliographic
  record and abstract were read; the scan did not expose
  machine-readable text in the browser. No claim that the whole paper
  was read. Its original reduction formulation needs care before
  attributing a modern many-one statement verbatim to it.
\item
  Richard Karp, \emph{Reducibility Among Combinatorial Problems} (1972):
  bibliographic identification only in this pass. Obtain and read the
  original relevant construction before expanding the full hardness
  chain.
\end{itemize}

\subsubsection{Original derivations and
computations}\label{original-derivations-and-computations}

\paragraph{Second-pass source check}\label{second-pass-source-check}

Toronto CS 2401 (Fall 2015), instructor Toniann Pitassi, lecturer Thomas
Watson, scribe Benett Axtell:
\href{https://www.cs.toronto.edu/~toni/Courses/Complexity2015/lectures/lecture3.pdf}{Lecture
3}, section 4 on physical page 4. Read the definition of coNP through
language complementation, its universal-certificate formulation, and the
statement that P is closed under complementation. Used only for these
class conventions. The manuscript supplies its own deterministic
complement-closure argument. The scribe prose about ``flip the bits of
the witness'' is imprecise and is not copied or used as a proof: it is
the verifier predicate that must be negated under the changed
quantifier. No interactive-proof theorem is invoked.

Rechecked Sipser lecture 14's representation and complement question.
The new figures show explicit counting units and quantifiers, without a
separation diagram. The edge-deletion figure is generated from the
actual retained-graph trace; its test replays every step against
independent enumeration on all 4,096 four-vertex directed graphs.

The four-vertex instance, subset recurrence proof, dense reachable-state
count, elementary vertex-position CNF clause count and edge-deletion
self-reduction explanation are derived explicitly in the manuscript and
checked by the local reference. These are pedagogical reconstructions of
standard ideas, not claims of novel algorithms.

\texttt{reference.py} uses only Python's standard library. Its oracle
enumerates interior permutations independently of the subset algorithm.
Tests cover all 4,096 simple directed four-vertex graphs for existence
agreement, all 64 three-vertex graphs with all three reduction pivots,
and all 512 Boolean assignments per three-vertex CNF instance. Counts
are algorithmic counters, not timing data.

\subsubsection{Scope exclusions}\label{scope-exclusions}

No P-versus-NP resolution, best-known exact-algorithm claim, universal
molecular lower bound or new laboratory result. The full reduction
chain, deeper work/depth treatment and physical accounting remain
publication dependencies. Earlier historical/chemical details continue
to rely on the preserved Chapter 2 source ledger; this draft does not
re-report that work as new research.
